\documentclass[UTF8,openany,12pt]{ctexbook}

\usepackage[a4paper,margin=2.6cm]{geometry}
\usepackage{amsmath,amssymb,amsthm,mathtools,bm}
\usepackage{booktabs,array}
\usepackage{enumitem}
\usepackage{tikz}
\usetikzlibrary{arrows.meta,positioning,calc,trees,shapes.geometric,fit,backgrounds}
\usepackage[ruled,vlined,linesnumbered]{algorithm2e}
\usepackage[numbers,sort&compress]{natbib}
\usepackage[colorlinks=true,linkcolor=blue!60!black,citecolor=blue!60!black,urlcolor=blue!60!black]{hyperref}
\usepackage[nameinlink,noabbrev]{cleveref}

\ctexset{
  chapter = {name = {第,章}},
  section = {format = \Large\bfseries}
}

\numberwithin{equation}{chapter}

\theoremstyle{definition}
\newtheorem{definition}{定义}[chapter]
\newtheorem{example}{例}[chapter]
\newtheorem{counterexample}{反例}[chapter]
\newtheorem{exercise}{练习}[chapter]

\theoremstyle{plain}
\newtheorem{theorem}{定理}[chapter]
\newtheorem{proposition}[theorem]{命题}
\newtheorem{lemma}[theorem]{引理}
\newtheorem{corollary}[theorem]{推论}

\crefformat{chapter}{第~#2#1#3~章}
\Crefformat{chapter}{第~#2#1#3~章}
\crefformat{section}{第~#2#1#3~节}
\Crefformat{section}{第~#2#1#3~节}
\crefname{definition}{定义}{定义}
\crefname{theorem}{定理}{定理}
\crefname{proposition}{命题}{命题}
\crefname{lemma}{引理}{引理}
\crefname{corollary}{推论}{推论}
\crefname{example}{例}{例}
\crefname{counterexample}{反例}{反例}
\crefname{equation}{式}{式}
\crefname{figure}{图}{图}
\crefname{table}{表}{表}
\crefname{algocf}{算法}{算法}

\SetKwInOut{KwIn}{输入}
\SetKwInOut{KwOut}{输出}
\SetKw{KwRet}{返回}
\SetKw{KwAnd}{且}
\SetKw{KwOr}{或}
\SetKwComment{tcp}{// }{}
\DontPrintSemicolon

\newcommand{\R}{\mathbb{R}}
\newcommand{\E}{\mathbb{E}}
\newcommand{\Var}{\operatorname{Var}}
\newcommand{\Cov}{\operatorname{Cov}}
\newcommand{\argmin}{\operatorname*{arg\,min}}
\newcommand{\argmax}{\operatorname*{arg\,max}}
\newcommand{\dist}{\operatorname{dist}}
\newcommand{\pa}{\operatorname{pa}}
\newcommand{\ch}{\operatorname{ch}}
\newcommand{\lca}{\operatorname{lca}}
\newcommand{\pathset}{\mathcal{P}}

\title{配电网拓扑辨识中的图论基础：从最小生成树到隐藏树重构}
\author{研究助理生成稿}
\date{\today}

\begin{document}

\frontmatter
\maketitle
\tableofcontents

\mainmatter
\chapter{绪论：为什么配电网拓扑辨识需要图论}
\label{chap:intro}

\section*{本章目标}
本章建立本书的共同语言：把配电网拓扑辨识表述为图学习问题，区分全节点可观测、仅叶节点可观测与候选图内开关状态辨识三类任务，并说明为什么后续的最小生成树、加性树距离和隐藏树重构不是抽象数学装饰，而是由放射状配电网的物理结构自然引出的工具。后文的树论基础见\cref{chap:trees}，最小生成树理论见\cref{chap:mst}。

\section{拓扑辨识问题的图论表述}
\label{sec:topology_problem}

配电网拓扑辨识的目标，是根据量测数据与先验信息恢复电气连接关系。把母线、配变、分支箱、智能电表等对象抽象为节点，把线路、开关、熔断器或等效馈线段抽象为边，就得到图
\begin{equation}
  G=(V,E),
\end{equation}
其中 $V$ 是节点集合，$E$ 是边集合。若网络在某个运行时段呈放射状，则实际带电拓扑通常是一棵以变电站或台区总表为根的树
\begin{equation}
  T=(V,E_T), \qquad |E_T|=|V|-1 .
\end{equation}

\begin{definition}[配电网拓扑辨识]
给定时序量测 $\mathcal{Y}$、可能的候选边集合 $E_0$ 以及先验信息 $\mathcal{I}$，拓扑辨识是估计实际电气边集
\begin{equation}
  \widehat E_T = \mathcal{A}(\mathcal{Y},E_0,\mathcal{I})
\end{equation}
的过程。若 $\widehat E_T=E_T$，称辨识精确；若二者在隐藏度为 $2$ 节点压缩意义下等价，则称辨识得到等效拓扑。
\end{definition}

这里的“拓扑”不是地理上的相邻，也不是通信链路上的可达，而是电气潮流方程所使用的节点导纳关系。这个区分贯穿全书：\cref{chap:multiview}将专门讨论通信图与电气图不相同的情形。

\section{放射状网、树与有根树}

低压和中压配电网常以放射状方式运行。图论中的树恰好刻画这种结构。

\begin{definition}[树和有根树]
无向连通且无环的图称为树。若指定一个节点 $0\in V$ 为根，则树成为有根树。对非根节点 $i$，从 $0$ 到 $i$ 的唯一路径上与 $i$ 相邻且更靠近根的节点称为 $i$ 的父节点，记为 $\pa(i)$；更远离根的相邻节点称为子节点，集合记为 $\ch(i)$。
\end{definition}

在配电网中，根节点通常对应变电站、馈线首端或台区总表；叶节点通常对应末端用户、电表箱或没有下游支路的节点；内部节点对应分支点、接线箱、开关柜或未装表的中间电气节点。隐藏内部节点是否可恢复，是本书从最小生成树转向隐藏树重构的核心原因，见\cref{chap:hidden_tree}。

\section{末端智能电表量测场景}

实际台区中，智能电表往往部署在用户端。设可观测节点集合为 $L\subseteq V$，通常 $L$ 接近叶节点集合，而内部分支节点 $H=V\setminus L$ 未量测。可得到的时序数据包括有功功率 $P_i(t)$、无功功率 $Q_i(t)$、电压幅值 $V_i(t)$、电流或功率因数等。对仅叶节点可观测场景，算法直接在 $L$ 上建树只会得到一棵连接叶节点的树，而不是原始电气树 $T$。因此，若真实网络含有隐藏分支点，必须估计或重构这些隐藏节点。

\begin{example}[一条馈线的两种观测图]
设真实拓扑为 $0$--$a$--$\{1,2\}$ 与 $a$--$b$--$\{3,4\}$，其中 $0$ 是根，$a,b$ 是分支节点，$1,2,3,4$ 是用户电表。若只有 $1,2,3,4$ 的电压与功率曲线，观测图节点集是 $\{1,2,3,4\}$。即使可以精确得到任意两个电表间的电气距离，也仍需判断 $1,2$ 是否共享隐藏分支点、$3,4$ 是否共享另一个隐藏分支点，这已经超出普通 MST 的输出形式。
\end{example}

\section{多源信息与图关系学习}

拓扑辨识中常见的信息源可以分为四类。

\begin{enumerate}[label=(\arabic*)]
  \item \textbf{电气量测}：$P/Q/V/I$ 时序、停复电事件、电压扰动响应。它们通过潮流方程与电气拓扑相关，是本书\cref{chap:impedance_distance,chap:estimation}的重点。
  \item \textbf{通信量测}：PLC/HPLC 的 RSSI、SNR、RTT、丢包率、路由 hop、载波信道响应等。它们反映通信信道质量，不等同于电气边，见\cref{chap:multiview}。
  \item \textbf{GIS 与台账}：电缆走向、杆塔坐标、箱变归属、开关台账。这些信息给出候选边和物理约束，但可能存在漏录、错录和时效性问题。
  \item \textbf{运行约束}：放射状约束、开关状态约束、潮流可行性和电压上下限约束。它们可用于后处理或联合优化。
\end{enumerate}

因此，一个现代拓扑辨识模型可被看成多视图图关系学习问题：
\begin{equation}
  p_{ij} = \mathbb{P}\{(i,j)\in E_T \mid
  \phi^{V}_{ij},\phi^{PQ}_{ij},\phi^{PLC}_{ij},\phi^{GIS}_{ij}\},
\end{equation}
再把边概率解码为满足放射状约束的图。若概率越大越可信，可用最大生成树；若距离越小越可信，可用最小生成树。

\section{三类拓扑辨识任务}

\subsection{全节点可观测}
若 $V$ 中每个电气节点均有量测，且可构造精确的加性树距离，那么在完全图 $K_V$ 上求 MST 可以恢复真实树。这个结论将在\cref{chap:mst_recovery}严格证明。该情形是 MST 方法最干净、最有理论保证的适用场景。

\subsection{仅叶节点可观测}
若只观测叶节点 $L$，距离矩阵定义在 $L\times L$ 上。真实拓扑含隐藏节点时，$K_L$ 上的 MST 至多产生一棵叶节点树，无法显式恢复隐藏分支点。此时应使用 four-point condition、additive phylogeny、neighbor joining 或 recursive grouping 等隐藏树重构思想，见\cref{chap:four_point,chap:hidden_tree}。

\subsection{候选图内开关状态辨识}
若 GIS 给出候选图 $G_0=(V,E_0)$，任务变为在候选边中选择一组闭合边 $E_T\subseteq E_0$。这类问题可加入放射状约束
\begin{equation}
  |E_T|=|V|-1,\qquad (V,E_T)\ \text{连通},
\end{equation}
也可用混合整数规划、最大生成树、最短路树或贝叶斯图模型求解。

\section{本章小结}

配电网拓扑辨识的数学核心是：用电气、通信和地理信息推断树形或近似树形的电气图。全节点可观测且距离是加性树距离时，MST 有清晰保证；仅叶节点可观测时，隐藏节点导致普通 MST 输出对象不匹配；候选图内开关状态辨识则更像受约束的图选择问题。后续章节将依次补齐树、MST、树距离、阻抗距离与多视图解码的理论基础。

\section*{思考题}
\begin{enumerate}
  \item 为什么“放射状运行”可以抽象为树？若存在联络开关但处于断开状态，它应如何进入图模型？
  \item 若台区总表与所有用户表均可观测，但中间分支箱不可观测，这属于全节点可观测还是叶节点可观测？
  \item 举一个通信图中相邻但电气图中不相邻的例子，并解释可能原因。
  \item 在候选图 $G_0$ 内辨识开关状态时，为什么只根据边概率逐边阈值化可能破坏放射状约束？
\end{enumerate}

\chapter{树、路径与有根树基础}
\label{chap:trees}

\section*{本章目标}
本章系统介绍图、树、路径、有根树和最近公共祖先等概念，并把它们翻译成配电网语言。读完本章后，应能判断一个隐藏分支节点是否具有可辨识意义，并理解\cref{chap:mst}中 MST 算法为什么能利用割和环的结构。

\section{图、路径、连通性与无环性}

\begin{definition}[无向图]
无向图 $G=(V,E)$ 由有限节点集 $V$ 和无向边集 $E\subseteq \{\{i,j\}:i,j\in V,i\ne j\}$ 构成。若 $e=\{i,j\}\in E$，称 $i$ 与 $j$ 相邻。
\end{definition}

\begin{definition}[路径和简单路径]
从 $v_0$ 到 $v_k$ 的路径是节点序列 $(v_0,v_1,\ldots,v_k)$，满足 $\{v_{\ell-1},v_\ell\}\in E$。若路径中节点互不重复，则称为简单路径。路径长度可以指边数，也可以指边权和；本书会根据上下文说明。
\end{definition}

\begin{definition}[连通与环]
若任意两个节点之间均存在路径，则图连通。若存在 $k\ge 3$ 且 $v_0=v_k$、$v_0,\ldots,v_{k-1}$ 互异的路径，则称图含环。
\end{definition}

\begin{theorem}[树的等价刻画]
\label{thm:tree_equiv}
对有限无向图 $T=(V,E)$，以下命题等价：
\begin{enumerate}[label=(\roman*)]
  \item $T$ 连通且无环；
  \item 任意两个节点之间存在唯一简单路径；
  \item $T$ 连通且 $|E|=|V|-1$；
  \item $T$ 无环且 $|E|=|V|-1$。
\end{enumerate}
\end{theorem}

\begin{proof}
若 $T$ 连通且无环，则任意两点至少有一条路径；若有两条不同简单路径，它们合并会形成环，矛盾，所以路径唯一。若任意两点有唯一简单路径，则图连通；若有环，环上两点之间有两条不同路径，矛盾，所以无环。连通无环图从一个节点开始逐步加入新节点，每加入一个节点必须且只需一条连接到已有部分的边，因此有 $|E|=|V|-1$。反过来，连通图若含环，删去环上一条边仍连通，会得到边数至少 $|V|$ 的冗余结构；无环图若不连通，则每个连通分量都是树，边数为 $\sum_r(|V_r|-1)=|V|-c<|V|-1$。于是得到等价性。
\end{proof}

\section{有根树、父子关系与深度}

指定根节点 $0$ 后，树中任意节点 $i$ 到根有唯一路径，记为
\begin{equation}
  \pathset_{0i}=(0=v_0,v_1,\ldots,v_k=i).
\end{equation}
节点 $i$ 的深度定义为 $\operatorname{depth}(i)=k$。若边带正权 $w_e$，也可定义加权深度
\begin{equation}
  h(i)=\sum_{e\in \pathset_{0i}} w_e .
\end{equation}
配电网中，深度常对应离电源的电气距离或线路层级；加权深度可对应公共路径阻抗和。

\begin{definition}[最近公共祖先]
在有根树中，节点 $i,j$ 的最近公共祖先 $\lca(i,j)$ 是同时位于 $\pathset_{0i}$ 与 $\pathset_{0j}$ 上、且深度最大的节点。
\end{definition}

\begin{proposition}[树上距离的 LCA 表达]
\label{prop:lca_distance}
若每条边权 $w_e>0$，令 $h(i)$ 为根到 $i$ 的加权深度，则
\begin{equation}
  d_T(i,j)=h(i)+h(j)-2h(\lca(i,j)).
\end{equation}
\end{proposition}

\begin{proof}
从 $i$ 到 $j$ 的唯一路径由 $i$ 上行到 $\lca(i,j)$，再下行到 $j$。第一段权重为 $h(i)-h(\lca(i,j))$，第二段为 $h(j)-h(\lca(i,j))$，相加即得结论。
\end{proof}

\section{配电网中的隐藏分支节点}

\begin{figure}[htbp]
  \centering
  \begin{tikzpicture}[
  node distance=1.35cm and 1.45cm,
  every node/.style={font=\small},
  obs/.style={circle,draw,fill=blue!18,minimum size=7mm},
  hid/.style={circle,draw,fill=white,minimum size=7mm},
  root/.style={circle,draw,fill=black!12,minimum size=7mm},
  line/.style={-Latex,thick}
]
\node[root] (r) {$0$};
\node[hid,below=of r] (a) {$a$};
\node[obs,below left=of a] (n1) {$1$};
\node[obs,below=of a] (n2) {$2$};
\node[hid,below right=of a] (b) {$b$};
\node[obs,below left=of b] (n3) {$3$};
\node[obs,below right=of b] (n4) {$4$};
\draw[thick] (r)--(a)--(n1);
\draw[thick] (a)--(n2);
\draw[thick] (a)--(b)--(n3);
\draw[thick] (b)--(n4);
\node[align=center,right=1.0cm of b] {隐藏分支点\\未直接量测};
\end{tikzpicture}

  \caption{配电网有根树与隐藏分支节点示意。实心节点表示可观测电表，空心节点表示未量测分支点。}
  \label{fig:distribution_tree_hidden}
\end{figure}

\Cref{fig:distribution_tree_hidden}展示了一个典型场景。根节点 $0$ 是台区总表，$a,b$ 是隐藏分支点，$1,2,3,4$ 是用户表。量测矩阵只包含用户侧曲线时，算法必须从叶节点间距离中恢复 $a,b$ 的存在。

\section{度为 2 隐藏节点为何不可唯一恢复}

\begin{definition}[度为 2 隐藏节点压缩]
若隐藏节点 $h$ 的度为 $2$，相邻边为 $(u,h)$ 与 $(h,v)$，且边权为 $w_{uh},w_{hv}$，则可把两条边压缩成一条边 $(u,v)$，权重为 $w_{uh}+w_{hv}$。压缩后所有可观测节点间的路径距离保持不变。
\end{definition}

\begin{proposition}[度为 2 隐藏节点不可由叶间距离唯一辨识]
\label{prop:degree2_unidentifiable}
若只观测某个集合 $L$ 上的两两路径距离，则任意不属于 $L$ 的度为 $2$ 隐藏节点及其相邻两条边的拆分比例不可唯一恢复。
\end{proposition}

\begin{proof}
设隐藏节点 $h$ 只连接 $u$ 与 $v$。对任意两个可观测节点 $i,j\in L$，若路径不经过 $h$，距离不变；若路径经过 $h$，贡献总是 $w_{uh}+w_{hv}$。把 $w_{uh},w_{hv}$ 替换为任意正数 $\alpha,\beta$ 且 $\alpha+\beta=w_{uh}+w_{hv}$，所有 $d(i,j)$ 完全相同。因此叶间距离不能识别 $h$ 的位置或两段边权。
\end{proof}

\begin{figure}[htbp]
  \centering
  \begin{tikzpicture}[
  every node/.style={font=\small},
  obs/.style={circle,draw,fill=blue!18,minimum size=7mm},
  hid/.style={circle,draw,fill=white,minimum size=7mm}
]
\node[obs] (u) at (0,0) {$u$};
\node[hid] (h) at (1.8,0) {$h$};
\node[obs] (v) at (3.6,0) {$v$};
\draw[thick] (u)--node[above] {$\alpha$} (h)--node[above] {$\beta$} (v);
\node at (4.8,0) {$\Longleftrightarrow$};
\node[obs] (u2) at (6.0,0) {$u$};
\node[obs] (v2) at (8.8,0) {$v$};
\draw[thick] (u2)--node[above] {$\alpha+\beta$} (v2);
\end{tikzpicture}

  \caption{度为 $2$ 隐藏节点压缩前后对端点距离等价。}
  \label{fig:degree2_hidden}
\end{figure}

\begin{example}[线路中点是否为节点]
若一条从分支箱到用户的电缆中间有一个未装表接头，且该接头不产生新的支路，则从拓扑辨识角度它是度为 $2$ 隐藏节点。仅靠端点电压功率曲线，不能判断这条线路是“一段”还是“两段”，只能估计总阻抗。
\end{example}

\section{本章小结}

树的关键特征是任意两点间路径唯一，这使路径距离、LCA 和公共路径阻抗都有明确含义。有根树把电源端方向引入图模型，便于解释父节点、子节点和分支层级。隐藏节点若度至少为 $3$，可能从叶间距离中体现为多个叶节点簇共享分支；若度为 $2$，则只能恢复压缩后的等效拓扑。这一事实解释了为什么\cref{chap:hidden_tree}讨论的隐藏树重构通常默认隐藏节点度至少为 $3$。

\section*{思考题}
\begin{enumerate}
  \item 证明：树中任意一条边都是割边，即删除它会使图不连通。
  \item 对\cref{fig:distribution_tree_hidden}，写出 $\lca(1,2)$、$\lca(1,4)$ 与 $\lca(3,4)$。
  \item 若一个隐藏节点度为 $4$，它是否一定可由叶间距离稳定恢复？还需要哪些条件？
  \item 解释为什么线路接头、熔断器、分支箱在不同数据粒度下可能对应不同的图节点定义。
\end{enumerate}

\chapter{最小生成树理论}
\label{chap:mst}

\section*{本章目标}
本章介绍生成树、最小生成树、最大生成树、Kruskal 算法、Prim 算法以及 cut property 和 cycle property。掌握这些结果后，\cref{chap:mst_recovery}中的拓扑恢复定理会变得非常自然。

\section{生成树与最小生成树}

\begin{definition}[生成树]
设 $G=(V,E)$ 是连通无向图。若 $T=(V,E_T)$ 是 $G$ 的子图且 $T$ 是树，则称 $T$ 是 $G$ 的生成树。
\end{definition}

给定边权函数 $c:E\to \R$，生成树 $T$ 的总权重为
\begin{equation}
  c(T)=\sum_{e\in E_T} c_e .
\end{equation}
最小生成树（minimum spanning tree, MST）是总权重最小的生成树：
\begin{equation}
  T^\star \in \argmin_{T\text{ 是 }G\text{ 的生成树}} c(T).
\end{equation}
若把目标改为最大化，则得到最大生成树。最大生成树常用于边概率解码：当 $c_{ij}$ 是边存在概率或相似度时，选择总相似度最大的放射状结构。

\begin{example}[距离与相似度的符号]
若 $c_{ij}$ 表示电气距离，边权越小越可信，应求 MST。若 $s_{ij}$ 表示“相邻概率”或相关相似度，边权越大越可信，应求最大生成树。最大生成树可等价地在权重 $c_{ij}=-s_{ij}$ 上求 MST。
\end{example}

\section{Kruskal 算法}

Kruskal 算法按边权从小到大扫描边，只加入不会形成环的边。

\begin{algorithm}[htbp]
  \caption{Kruskal 最小生成树算法}
  \label{alg:kruskal}
  \KwIn{连通无向图 $G=(V,E)$，边权 $c_e$}
  \KwOut{一棵 MST $T=(V,E_T)$}
  $E_T\leftarrow \varnothing$\;
  初始化并查集，使每个节点自成一个连通分量\;
  将 $E$ 按 $c_e$ 非降序排序\;
  \ForEach{$e=(u,v)\in E$ 按排序顺序}{
    \If{$u$ 与 $v$ 不在同一并查集分量}{
      $E_T\leftarrow E_T\cup\{e\}$\;
      合并 $u$ 与 $v$ 的分量\;
    }
    \If{$|E_T|=|V|-1$}{停止}
  }
  \KwRet{$T=(V,E_T)$}
\end{algorithm}

若使用排序和近似常数时间并查集，复杂度为
\begin{equation}
  O(|E|\log |E|).
\end{equation}
在完全图 $K_V$ 上，$|E|=O(|V|^2)$，排序可能成为主要开销。

\section{Prim 算法}

Prim 算法从一个初始节点出发，每次选择连接当前树与外部节点的最小边。

\begin{algorithm}[htbp]
  \caption{Prim 最小生成树算法}
  \label{alg:prim}
  \KwIn{连通无向图 $G=(V,E)$，边权 $c_e$，起点 $r$}
  \KwOut{一棵 MST $T=(V,E_T)$}
  $S\leftarrow\{r\}$，$E_T\leftarrow\varnothing$\;
  \While{$S\ne V$}{
    选择最小权跨割边 $e^\star=(u,v)$，其中 $u\in S,v\notin S$\;
    $E_T\leftarrow E_T\cup\{e^\star\}$，$S\leftarrow S\cup\{v\}$\;
  }
  \KwRet{$T=(V,E_T)$}
\end{algorithm}

使用二叉堆实现时，复杂度可达 $O(|E|\log |V|)$；对稠密图也可用邻接矩阵实现为 $O(|V|^2)$。在配电网拓扑恢复中，若距离矩阵已给出，完全图上 Prim 的矩阵实现很直接。

\section{割性质与环性质}

\begin{definition}[割与跨割边]
对非空真子集 $S\subset V$，割 $(S,V\setminus S)$ 的跨割边集合为
\begin{equation}
  \delta(S)=\{(u,v)\in E:u\in S,\ v\in V\setminus S\}.
\end{equation}
\end{definition}

\begin{theorem}[Cut property]
\label{thm:cut_property}
给定连通无向加权图 $G$。若边 $e$ 是某个割 $(S,V\setminus S)$ 上的唯一最小权跨割边，则 $e$ 属于每一棵 MST。
\end{theorem}

\begin{proof}
任取一棵 MST $T$。若 $e\in T$，结论成立。若 $e\notin T$，把 $e$ 加入 $T$ 会形成唯一环，该环必须还包含至少一条跨越同一割的边 $f$，否则从 $S$ 出发沿环经过 $e$ 后无法回到 $S$。由于 $e$ 是唯一最小跨割边，有 $c_e<c_f$。令
\begin{equation}
  T'=T\cup\{e\}\setminus\{f\}.
\end{equation}
则 $T'$ 仍为生成树，且 $c(T')=c(T)+c_e-c_f<c(T)$，与 $T$ 是 MST 矛盾。因此 $e\in T$。
\end{proof}

\begin{theorem}[Cycle property]
\label{thm:cycle_property}
若边 $e$ 是某个环上的唯一最大权边，则 $e$ 不属于任何 MST。
\end{theorem}

\begin{proof}
若某棵 MST $T$ 含 $e$，删除 $e$ 后 $T$ 分成两个连通分量。原环上除 $e$ 外必有一条边 $f$ 跨越这两个分量，否则环无法闭合。因为 $e$ 是唯一最大权边，$c_f<c_e$。用 $f$ 替换 $e$ 得到权重更小的生成树，矛盾。
\end{proof}

\section{唯一 MST 条件}

\begin{theorem}[互异边权推出 MST 唯一]
\label{thm:distinct_unique_mst}
若连通无向图 $G$ 的所有边权两两互异，则 MST 唯一。
\end{theorem}

\begin{proof}
反设存在两棵不同的 MST，记为 $T_1,T_2$。取对称差 $T_1\triangle T_2$ 中权重最小的边 $e$，不妨设 $e\in T_1\setminus T_2$。把 $e$ 加入 $T_2$ 得到唯一环 $C$。环 $C$ 中至少有一条边 $f\in T_2\setminus T_1$，否则 $T_1$ 本身含环。由于 $e$ 是对称差中最小边且边权互异，有 $c_e<c_f$。于是 $T_2\cup\{e\}\setminus\{f\}$ 是权重更小的生成树，矛盾。因此 MST 唯一。
\end{proof}

\section{一个配电网拓扑恢复例子}

\begin{figure}[htbp]
  \centering
  \begin{tikzpicture}[
  every node/.style={font=\small},
  bus/.style={circle,draw,fill=blue!12,minimum size=7mm}
]
\node[bus] (n0) at (0,0) {$0$};
\node[bus] (n1) at (2,0) {$1$};
\node[bus] (n2) at (4,0.8) {$2$};
\node[bus] (n3) at (4,-0.8) {$3$};
\draw[very thick] (n0)--node[above] {$1$} (n1);
\draw[very thick] (n1)--node[above] {$2$} (n2);
\draw[very thick] (n1)--node[below] {$3$} (n3);
\draw[dashed,gray] (n0) to[bend left=18] node[above,gray] {$3$} (n2);
\draw[dashed,gray] (n0) to[bend right=18] node[below,gray] {$4$} (n3);
\draw[dashed,gray] (n2)--node[right,gray] {$5$} (n3);
\end{tikzpicture}

  \caption{真实树距离上的 MST 恢复例子。完全图中未画出的边权为树上路径和，真实边是其基本割上的唯一最小边。}
  \label{fig:mst_example}
\end{figure}

\begin{example}[四节点馈线]
设真实配电树为 $0$--$1$--$2$ 和 $1$--$3$，边权分别为 $1,2,3$。由树路径和得到
\[
  d_{01}=1,\quad d_{12}=2,\quad d_{13}=3,\quad
  d_{02}=3,\quad d_{03}=4,\quad d_{23}=5.
\]
在完全图 $K_{\{0,1,2,3\}}$ 上以 $d_{ij}$ 为边权求 MST，Kruskal 首先选择 $(0,1)$，再选 $(1,2)$，再选 $(1,3)$，正好恢复真实树。这个例子是\cref{thm:additive_mst_recovery}的最小版本。
\end{example}

\section{本章小结}

MST 是在连通图上选择总权最小树的经典问题。Kruskal 从小边开始避免成环，Prim 从一个连通块开始反复跨割扩张。Cut property 说明唯一最小跨割边必进所有 MST；cycle property 说明唯一最大环边必不进 MST。边权互异是 MST 唯一的充分条件，但不是必要条件。下一章将利用 cut property 证明：若完整节点集上的距离是精确加性树距离，则 MST 可以恢复真实配电树。

\section*{思考题}
\begin{enumerate}
  \item 在一个三角形图中，边权分别为 $1,1,2$。MST 是否唯一？为什么？
  \item Kruskal 和 Prim 的输出是否一定相同？在什么条件下相同？
  \item 用 cycle property 解释为什么在距离图上很长的跨支路边不应进入 MST。
  \item 若边权是估计误差较大的电气距离，MST 会优先受哪些边影响？
\end{enumerate}

\chapter{MST 何时能恢复真实配电网拓扑}
\label{chap:mst_recovery}

\section*{本章目标}
本章给出 MST 恢复真实配电网树的严格条件。核心结论是：当完整节点集上的边权来自真实树上的精确加性路径距离时，在完全图上求 MST 会恢复原树；当距离带噪时，恢复稳定性由每条真实边的割间隔控制。本章结果建立在\cref{chap:mst}的 cut property 上，并为\cref{chap:counterexamples}中的失败情形提供对照。

\section{完整节点集与加性树距离}

设真实拓扑为树 $T=(V,E_T)$，每条真实线路 $e\in E_T$ 带正权 $w_e>0$。树上路径距离定义为
\begin{equation}
  d_{ij}=\sum_{e\in \pathset_{ij}} w_e,\qquad i,j\in V,
  \label{eq:additive_distance}
\end{equation}
其中 $\pathset_{ij}$ 是 $T$ 中从 $i$ 到 $j$ 的唯一路径。若配电网中 $w_e$ 表示线路电阻、电抗或某种可加阻抗，则 $d_{ij}$ 是两节点之间路径阻抗和。

注意这里的节点集是完整的：每个真实节点，包括中间分支节点，都在 $V$ 中并可进入距离矩阵。若只观测叶节点，结论不能直接使用，见\cref{chap:hidden_tree}。

\begin{theorem}[完整节点集上的精确加性树距离可由 MST 恢复真实树]
\label{thm:additive_mst_recovery}
设真实拓扑 $T=(V,E_T)$ 为一棵树，每条边权 $w_e>0$。若对任意 $i,j\in V$，距离 $d_{ij}$ 等于 $T$ 中 $i$ 到 $j$ 路径上的边权和，则在完全图 $K_V$ 上以 $d_{ij}$ 为边权求 MST，所得生成树等于 $T$。
\end{theorem}

\begin{proof}
任取真实边 $e=(u,v)\in E_T$。删除 $e$ 后，树 $T$ 被分成两个连通分量，记为 $A_e$ 与 $B_e$，其中 $u\in A_e$、$v\in B_e$。这给出完全图 $K_V$ 上的一个割 $(A_e,B_e)$。

对任意跨割边 $(a,b)$，其中 $a\in A_e$、$b\in B_e$，树中从 $a$ 到 $b$ 的唯一路径必须先从 $a$ 走到 $u$，经过真实边 $e=(u,v)$，再从 $v$ 走到 $b$。因此
\begin{equation}
  d_{ab}=d_{au}+w_e+d_{vb}.
  \label{eq:cut_decomposition}
\end{equation}
当 $(a,b)=(u,v)$ 时，$d_{ab}=w_e$。当 $(a,b)\ne (u,v)$ 时，至少有 $a\ne u$ 或 $b\ne v$。由于所有边权严格为正，若 $a\ne u$ 则 $d_{au}>0$，若 $b\ne v$ 则 $d_{vb}>0$，故
\begin{equation}
  d_{ab}=d_{au}+w_e+d_{vb}>w_e=d_{uv}.
\end{equation}
也就是说，真实边 $e=(u,v)$ 是割 $(A_e,B_e)$ 上唯一最小的跨割边。

由\cref{thm:cut_property}，唯一最小跨割边属于每一棵 MST。因此每条真实边 $e\in E_T$ 都属于完全图 $K_V$ 的任意 MST。真实树共有 $|V|-1$ 条边，而任意生成树也只有 $|V|-1$ 条边，所以 MST 的边集只能等于 $E_T$。故 MST 恢复真实树 $T$。
\end{proof}

\section{Cut-consistency 充分条件}

定理~\ref{thm:additive_mst_recovery}利用加性距离推出了“每条真实边在其基本割上唯一最小”。事实上，MST 恢复并不要求全局距离严格等于树路径和，只需要这个割一致性条件成立。

\begin{theorem}[MST 恢复的 cut-consistency 充分条件]
\label{thm:cut_consistency}
给定真实树 $T=(V,E_T)$ 与完全图 $K_V$ 上的边权 $c_{ij}$。对任意真实边 $e=(u,v)$，删除 $e$ 得到基本割 $(A_e,B_e)$。若 $e$ 都是其基本割上的唯一最小跨割边，即
\begin{equation}
  c_e < c_{ab},\qquad
  \forall a\in A_e,\ b\in B_e,\ (a,b)\ne e,
  \label{eq:cut_consistency}
\end{equation}
则在 $K_V$ 上以 $c_{ij}$ 为边权求 MST 会恢复真实树 $T$。
\end{theorem}

\begin{proof}
对每条真实边 $e\in E_T$，条件\eqref{eq:cut_consistency}说明 $e$ 是割 $(A_e,B_e)$ 上唯一最小跨割边。由 cut property，$e$ 属于所有 MST。于是所有 $|V|-1$ 条真实边均属于任意 MST。任意生成树只有 $|V|-1$ 条边，所以 MST 的边集必为 $E_T$。
\end{proof}

\begin{example}[非加性但割一致的权重]
真实树为链 $1$--$2$--$3$，令 $c_{12}=1$、$c_{23}=1$、$c_{13}=10$。这组权重是某棵树距离，但若把 $c_{13}$ 改成 $3$、$4$ 或 $100$，只要仍大于 $1$，MST 都会恢复真实链。MST 只关心每个基本割上的相对大小，而不是所有三角关系或 four-point condition。
\end{example}

\section{噪声下的 margin 条件}

实际距离矩阵来自回归、相关性或通信特征，通常含有误差。定义每条真实边 $e=(u,v)$ 的割间隔
\begin{equation}
  \gamma_e =
  \min_{\substack{a\in A_e,\ b\in B_e\\(a,b)\ne e}}
  \left(d_{ab}-d_e\right),
  \label{eq:cut_margin}
\end{equation}
其中 $(A_e,B_e)$ 是删除 $e$ 后的基本割。若 $d$ 是正边权加性树距离，则由\cref{thm:additive_mst_recovery}的证明可知 $\gamma_e>0$。最小间隔
\begin{equation}
  \gamma_{\min}=\min_{e\in E_T}\gamma_e
\end{equation}
越大，MST 对距离估计误差越稳健。

\begin{theorem}[带噪距离矩阵下 MST 恢复的 margin 条件]
\label{thm:noisy_margin_mst}
设 $d_{ij}$ 是完整节点集上的真实距离，并满足\cref{thm:cut_consistency}。对每条真实边按式\eqref{eq:cut_margin}定义 $\gamma_e$，且 $\gamma_{\min}=\min_e\gamma_e>0$。若估计距离 $\widehat d_{ij}$ 满足
\begin{equation}
  \max_{i,j\in V}|\widehat d_{ij}-d_{ij}|
  < \frac{1}{2}\gamma_{\min},
  \label{eq:noise_margin_bound}
\end{equation}
则在完全图 $K_V$ 上以 $\widehat d_{ij}$ 为边权求 MST 仍恢复真实树 $T$。
\end{theorem}

\begin{proof}
令
\begin{equation}
  \varepsilon=\max_{i,j\in V}|\widehat d_{ij}-d_{ij}|.
\end{equation}
任取真实边 $e=(u,v)$ 及其基本割 $(A_e,B_e)$。对任意跨割边 $(a,b)\ne e$，有
\begin{align}
  \widehat d_{ab}-\widehat d_e
  &= (d_{ab}-d_e)
     +(\widehat d_{ab}-d_{ab})
     -(\widehat d_e-d_e)  \notag\\
  &\ge \gamma_e - |\widehat d_{ab}-d_{ab}|-|\widehat d_e-d_e| \notag\\
  &\ge \gamma_{\min}-2\varepsilon .
\end{align}
由条件\eqref{eq:noise_margin_bound}，$\gamma_{\min}-2\varepsilon>0$，故
\begin{equation}
  \widehat d_e < \widehat d_{ab},\qquad
  \forall a\in A_e,\ b\in B_e,\ (a,b)\ne e .
\end{equation}
因此每条真实边在估计距离下仍是其基本割上的唯一最小跨割边。应用\cref{thm:cut_consistency}，MST 恢复真实树。
\end{proof}

\begin{corollary}[加性距离的可恢复噪声半径]
\label{cor:additive_noise_radius}
若 $d$ 来自正权真实树的加性路径距离，且估计误差满足式\eqref{eq:noise_margin_bound}，则估计距离上的 MST 恢复真实树。
\end{corollary}

\begin{proof}
由\cref{thm:additive_mst_recovery}，加性路径距离满足 cut-consistency，且 $\gamma_{\min}>0$。再由\cref{thm:noisy_margin_mst}即得。
\end{proof}

\section{配电网物理解释}

在配电网中，一条真实线路 $e=(u,v)$ 把网络分成上游侧 $A_e$ 与下游侧 $B_e$。若边权是线路电阻，跨越该割的任意两个节点 $a\in A_e$、$b\in B_e$ 的路径都必须经过线路 $e$，并额外包含 $a$ 到 $u$ 或 $v$ 到 $b$ 的路径。因此真实端点 $(u,v)$ 是跨割节点对中电气距离最短的一对。这就是 MST 能恢复拓扑的物理原因：它不是“相关性越大越相邻”的经验规则，而是“每条真实线路是其上下游割上的最近跨割连接”的结构事实。

噪声定理则说明，估计误差并不需要为零；只要误差小于最薄弱割间隔的一半，所有割上的最近跨割关系仍不改变。实际算法中，可把小间隔边视为脆弱边：短支线、阻抗相近的并行候选路径、未量测负荷严重的区域，往往具有较小 margin。

\section{MST 使用检查清单}

\begin{algorithm}[htbp]
  \caption{MST 拓扑恢复的理论检查流程}
  \label{alg:mst_checklist}
  \KwIn{节点集合 $V$，估计距离 $\widehat D$，可用先验 $\mathcal{I}$}
  \KwOut{是否适合直接使用 MST 的判断}
  检查 $V$ 是否接近完整电气节点集\;
  检查距离是否应为树上路径可加量，如电阻距离或电抗距离\;
  若有候选真实边或台账，估计基本割上的最小间隔\;
  若仅叶节点可观测，则转向隐藏树重构而非直接解释 MST 为原始拓扑\;
  若距离由相关性得到，先剔除公共模态并做物理一致性验证\;
  \KwRet{满足完整节点、加性距离、足够 margin 时，直接 MST 较可靠}
\end{algorithm}

\section{本章小结}

本章证明了三个核心结果。第一，完整节点集上的精确加性树距离保证 MST 恢复真实树。第二，更一般地，只要每条真实边是其基本割上的唯一最小跨割边，MST 就恢复真实树。第三，在距离带噪时，若最大误差小于最小割间隔的一半，MST 的输出仍不变。这些定理也指出了 MST 的边界：缺失隐藏节点、距离非加性、公共模态污染或强相关负荷都可能破坏割一致性。

\section*{思考题}
\begin{enumerate}
  \item 对链 $1$--$2$--$3$--$4$，边权均为 $1$，计算每条真实边的 $\gamma_e$。
  \item 若某条真实边的割间隔很小，实际数据采集应如何改进以提高辨识稳定性？
  \item 定理~\ref{thm:cut_consistency}是充分条件。请尝试构造一个 MST 正确但某条真实边不是基本割唯一最小边的例子，并说明为什么不矛盾。
  \item 在候选图不是完全图而是 GIS 图时，本章证明需要怎样修改？
\end{enumerate}

\chapter{反例：为什么“有距离矩阵”不代表 MST 一定正确}
\label{chap:counterexamples}

\section*{本章目标}
本章用五个反例说明 MST 的适用条件不可省略。与\cref{chap:mst_recovery}相比，本章关注条件失效后的具体后果：节点集不完整、距离非加性、公共源端电压模态、负荷高度相关以及通信图与电气图不一致。

\section{反例 1：仅叶节点可观测的星形隐藏中心}

\begin{counterexample}[隐藏中心星形结构]
\label{cex:leaf_star}
真实树为一个隐藏中心 $h$ 连接三个可观测叶节点 $1,2,3$，边权均为 $1$。叶间距离为
\begin{equation}
  d_{12}=d_{13}=d_{23}=2.
\end{equation}
在完全图 $K_{\{1,2,3\}}$ 上求 MST，只能得到三种等价输出之一，例如边集 $\{(1,2),(1,3)\}$。该输出把 $1$ 误当作中心节点，且没有恢复隐藏节点 $h$。
\end{counterexample}

这个反例不是算法实现问题，而是输出对象不匹配：MST 的节点集被固定为可观测叶节点，不能凭空产生隐藏分支点。正确做法是重构带隐藏节点的等效树，见\cref{chap:hidden_tree}。

\section{反例 2：距离不是加性树距离}

\begin{counterexample}[非加性权重导致选错边]
\label{cex:nonadditive}
真实树为 $1$--$2$--$3$，真实边权均为 $1$，因此加性距离应满足 $d_{13}=2$。假设某种经验特征给出
\begin{equation}
  \widehat d_{12}=1,\qquad \widehat d_{23}=1,\qquad \widehat d_{13}=0.2.
\end{equation}
Kruskal 会先选 $(1,3)$，再选 $(1,2)$ 或 $(2,3)$，得到的树不是真实树。
\end{counterexample}

非加性可能来自错误归一化、把相似度当距离、相关性未取负号、或使用了与路径阻抗无关的统计量。MST 只保证在输入权重满足割一致性时可靠；“任意距离矩阵”没有这种保证。

\section{反例 3：公共源端电压导致错误连接}

\begin{counterexample}[公共模态污染相关性]
\label{cex:source_voltage}
设两个末端节点 $i,j$ 分属不同支路，但都受到源端电压波动 $s(t)$ 的强影响：
\begin{equation}
  V_i(t)=s(t)+\eta_i(t),\qquad
  V_j(t)=s(t)+\eta_j(t),
\end{equation}
其中 $\eta_i,\eta_j$ 是较小的局部扰动。即使 $i$ 与 $j$ 在电气拓扑上相距较远，二者电压曲线相关系数也会很高。若用
\begin{equation}
  \widehat d_{ij}=1-\operatorname{corr}(V_i,V_j)
\end{equation}
作为距离，MST 可能优先连接这些共享源端模态但并不相邻的节点。
\end{counterexample}

物理上，源端电压波动是全局共同因子。拓扑辨识应使用电压增量对负荷增量的响应、公共模态去除或条件相关，而不能直接把原始电压相关性解释为线路相邻性。

\section{反例 4：负荷高度相关导致距离退化}

\begin{counterexample}[同步负荷曲线]
\label{cex:correlated_load}
设两个不同支路上的用户具有几乎同步的有功和无功变化：
\begin{equation}
  \Delta P_i(t)\approx \alpha_i z(t),\qquad
  \Delta P_j(t)\approx \alpha_j z(t).
\end{equation}
如果距离估计依赖 $\Delta V$ 与 $\Delta P$ 的协方差或回归，而设计矩阵列高度共线，则多个拓扑模型会解释出相近的电压变化。估计出的 $\widehat d_{ij}$ 可能失去区分度，甚至把不同支路用户聚为一簇。
\end{counterexample}

这个问题对应统计学中的病态回归和遗漏变量偏差。岭回归可以缓解方差但会引入偏差；Lasso 可能在相关特征中任意选择一个。第\cref{chap:estimation}将讨论这些问题如何影响 MST margin。

\section{反例 5：通信路由图不等于电气拓扑图}

\begin{counterexample}[PLC 路由边不是电气边]
\label{cex:communication_not_electrical}
在 HPLC 网络中，电表 $i$ 可能通过通信质量更好的中继电表 $j$ 转发数据，即路由表中存在 $i\to j$。但 $i$ 与 $j$ 的电气连接可能分别位于不同分支，只是载波信道在某个频段上更稳定。若把通信路由边直接当成电气边，得到的图可能含有跨支路边，也可能不满足放射状电气结构。
\end{counterexample}

通信特征有价值，但更适合作为 noisy prior，而不是硬等式。实际模型应学习
\begin{equation}
  \mathbb{P}\{(i,j)\in E_E\mid \text{电气特征},\text{通信特征},\text{GIS 特征}\},
\end{equation}
再由生成树约束解码，见\cref{chap:multiview}。

\section{一个统一视角：MST 失败的根源}

上述反例可归纳为三类失配。
\begin{enumerate}[label=(\arabic*)]
  \item \textbf{节点集失配}：真实树含隐藏内部节点，但算法只允许叶节点成树。
  \item \textbf{权重语义失配}：输入不是树路径可加距离，也不满足 cut-consistency。
  \item \textbf{数据生成机制失配}：公共模态、相关负荷、通信中继等因素改变了统计相似度与电气相邻性的关系。
\end{enumerate}

因此，正确的研究问题不是“是否能用 MST”，而是“在哪个节点集上、用何种物理距离、在多大估计误差下使用 MST”。这也是\cref{chap:impedance_distance,chap:estimation}要回答的问题。

\section{本章小结}

MST 是强有力的图算法，但它不会自动把任意相似度矩阵变成真实拓扑。若完整节点集和加性距离条件成立，MST 有严格保证；若只有叶节点、距离非加性或数据受公共模态污染，则应改用隐藏树重构、物理回归、公共因子校正或多视图约束模型。

\section*{思考题}
\begin{enumerate}
  \item 对\cref{cex:leaf_star}，如果把隐藏中心 $h$ 加入节点集，MST 会怎样输出？
  \item 构造一个四节点非加性距离矩阵，使 MST 输出一条跨支路错误边。
  \item 为什么“电压相关性高”更可能表示共同源端扰动，而不一定表示两节点直接相连？
  \item 在 PLC 场景中，哪些通信特征更可能与电气距离相关？哪些更容易受非电气因素影响？
\end{enumerate}

\chapter{加性树距离与 four-point condition}
\label{chap:four_point}

\section*{本章目标}
本章介绍 additive tree metric 和 four-point condition，说明为什么叶节点距离可以包含隐藏树结构信息，也说明为什么仅在叶节点上求 MST 不足以恢复隐藏分支节点。本章为\cref{chap:hidden_tree}的重构算法做准备。

\section{加性树距离}

\begin{definition}[加性树度量]
设 $T=(V,E)$ 是一棵树，边权 $w_e>0$。若有限集合 $L\subseteq V$ 上的距离 $d:L\times L\to \R_{\ge 0}$ 满足
\begin{equation}
  d_{ij}=\sum_{e\in \pathset_{ij}}w_e,\qquad i,j\in L,
\end{equation}
则称 $d$ 是由树 $T$ 在 $L$ 上诱导的加性树度量。若 $L=V$，称为完整树度量；若 $L$ 只含叶节点，称为叶间树度量。
\end{definition}

加性树度量的关键是“路径可加”。阻抗距离、线路长度距离和某些线性化潮流响应距离都可在合适条件下具有这种结构；普通欧氏坐标距离、通信 hop 距离或原始相关性距离则未必满足。

\section{Four-point condition}

\begin{definition}[Four-point condition]
距离矩阵 $d$ 满足 four-point condition，是指对任意四个不同点 $a,b,c,d$，三项
\begin{equation}
  S_1=d_{ab}+d_{cd},\qquad
  S_2=d_{ac}+d_{bd},\qquad
  S_3=d_{ad}+d_{bc}
\end{equation}
中的最大值至少出现两次。若四点形成非退化 quartet，则通常有两项最大且相等，另一项较小。
\end{definition}

\begin{figure}[htbp]
  \centering
  \begin{tikzpicture}[
  every node/.style={font=\small},
  leaf/.style={circle,draw,fill=blue!14,minimum size=7mm},
  hid/.style={circle,draw,fill=white,minimum size=6mm}
]
\node[hid] (u) at (0,0) {};
\node[hid] (v) at (2.2,0) {};
\node[leaf] (a) at (-1.1,0.9) {$a$};
\node[leaf] (b) at (-1.1,-0.9) {$b$};
\node[leaf] (c) at (3.3,0.9) {$c$};
\node[leaf] (d) at (3.3,-0.9) {$d$};
\draw[thick] (a)--node[above left] {$\alpha$} (u);
\draw[thick] (b)--node[below left] {$\beta$} (u);
\draw[very thick] (u)--node[above] {$x$} (v);
\draw[thick] (v)--node[above right] {$\gamma$} (c);
\draw[thick] (v)--node[below right] {$\delta$} (d);
\node at (1.1,-0.65) {$ab|cd$};
\end{tikzpicture}

  \caption{quartet 树 $ab|cd$：叶节点 $a,b$ 在内部边一侧，$c,d$ 在另一侧。}
  \label{fig:quartet}
\end{figure}

\begin{lemma}[树度量满足 four-point condition]
\label{lem:four_point_necessary}
任意加性树度量都满足 four-point condition。
\end{lemma}

\begin{proof}
任取四个叶节点 $a,b,c,d$。若它们在树上的最小连通子树没有内部边，即呈星形退化，则三项 $S_1,S_2,S_3$ 相等。否则该最小子树含一条内部边，诱导出一个 quartet 分裂。以\cref{fig:quartet}的 $ab|cd$ 为例，设内部边长度为 $x>0$，四条 pendant 边长度分别为 $\alpha,\beta,\gamma,\delta$。则
\begin{align}
  d_{ab}+d_{cd} &= (\alpha+\beta)+(\gamma+\delta),\\
  d_{ac}+d_{bd} &= (\alpha+x+\gamma)+(\beta+x+\delta)
                 = \alpha+\beta+\gamma+\delta+2x,\\
  d_{ad}+d_{bc} &= (\alpha+x+\delta)+(\beta+x+\gamma)
                 = \alpha+\beta+\gamma+\delta+2x.
\end{align}
因此后两项相等且大于第一项。其他分裂情形只是重命名节点，结论相同。
\end{proof}

\begin{proposition}[quartet 分裂的判别]
\label{prop:quartet_split}
若四点树度量中
\begin{equation}
  d_{ab}+d_{cd} < d_{ac}+d_{bd}=d_{ad}+d_{bc},
\end{equation}
则四点的非退化分裂为 $ab|cd$，内部边长度为
\begin{equation}
  x=\frac{1}{2}\left(d_{ac}+d_{bd}-d_{ab}-d_{cd}\right).
\end{equation}
\end{proposition}

\begin{proof}
由\cref{lem:four_point_necessary}证明中的 quartet 表达式，较小的和对应同侧叶对，较大两项比它多出 $2x$，故内部边长度如上。
\end{proof}

\section{与隐藏树重构的关系}

Four-point condition 的意义在于：即使内部节点不可观测，叶节点间距离仍然包含“哪些叶子共享近端分支”的信息。例如 $ab|cd$ 表示 $a,b$ 的路径在进入内部边前先汇合，$c,d$ 也先汇合。通过多个 quartet 的一致性，可以拼接出整棵隐藏树。

这解释了 leaf-only 场景不能只用 MST 的原因。MST 在叶节点集合上输出的是一棵叶节点树；hidden tree reconstruction 则允许新建内部节点，使输出树的叶间距离匹配输入距离。二者优化对象不同：
\begin{align}
  \text{MST:}\quad &\min_{T\text{ spans }L}\sum_{e\in T}\widehat d_e,\\
  \text{hidden tree:}\quad &\text{find }(\widetilde T,\widetilde w)\text{ such that }
  d_{\widetilde T}(i,j)\approx \widehat d_{ij},\ i,j\in L.
\end{align}

\section{四叶节点数值示例}

\begin{example}[判断 quartet]
设叶节点距离为
\[
\begin{array}{c|cccc}
 & a&b&c&d\\ \hline
a&0&2&5&5\\
b&2&0&5&5\\
c&5&5&0&2\\
d&5&5&2&0
\end{array}
\]
计算
\[
S_1=d_{ab}+d_{cd}=4,\quad
S_2=d_{ac}+d_{bd}=10,\quad
S_3=d_{ad}+d_{bc}=10.
\]
最大值出现两次，且 $S_1$ 最小，因此分裂为 $ab|cd$。内部边长度为 $(10-4)/2=3$。若四条 pendant 边均为 $1$，内部边为 $3$，正好得到该距离矩阵。
\end{example}

\begin{exercise}[给定矩阵判断 four-point condition]
对距离矩阵
\[
\begin{array}{c|cccc}
 &1&2&3&4\\ \hline
1&0&4&7&7\\
2&4&0&7&7\\
3&7&7&0&2\\
4&7&7&2&0
\end{array}
\]
判断是否满足 four-point condition，并给出 quartet 分裂与内部边长度。
\end{exercise}

\section{本章小结}

加性树度量把树结构编码进两两距离。Four-point condition 是树度量的基本判据：对任意四点，三种配对距离和的最大值至少出现两次；非退化时，较小项揭示 quartet 分裂。该条件使我们能够从叶节点距离中推断隐藏内部结构，也解释了为什么仅在叶节点上求 MST 无法恢复隐藏分支点。

\section*{思考题}
\begin{enumerate}
  \item 证明：若四个叶节点构成星形树，则 $S_1=S_2=S_3$。
  \item 若观测距离含噪，$S_2$ 与 $S_3$ 不再完全相等，应如何设置容差判断 quartet？
  \item Four-point condition 是树度量的必要条件。查阅 Buneman 相关结果，思考它在有限度量上是否也是充分条件。
  \item 在配电网中，哪些物理因素可能使叶间阻抗距离偏离精确树度量？
\end{enumerate}

\chapter{叶节点距离下的隐藏树重构}
\label{chap:hidden_tree}

\section*{本章目标}
本章讨论仅末端量测时如何从叶节点距离恢复隐藏分支结构。我们介绍 terminal nodes、Steiner/hidden nodes、additive phylogeny、neighbor joining 与 recursive grouping 的基本思想，并给出一个小规模数值例子。与\cref{chap:mst_recovery}不同，本章的输出允许包含未观测内部节点。

\section{Terminal nodes 与 hidden nodes}

设可观测节点集合为 $L$，称为 terminal nodes；未量测但可能影响叶间距离结构的内部节点集合为 $H$，称为 hidden nodes 或 Steiner nodes。隐藏树重构的目标是寻找
\begin{equation}
  \widetilde T=(L\cup \widetilde H,\widetilde E),\qquad \widetilde w_e>0,
\end{equation}
使得对任意 $i,j\in L$，
\begin{equation}
  d_{\widetilde T}(i,j)\approx \widehat d_{ij}.
\end{equation}

\begin{definition}[等效拓扑]
给定只在 terminal 集合 $L$ 上观测的距离，若两棵加权树在 $L$ 上诱导完全相同的两两距离，则称它们对 $L$ 等效。通常把所有度为 $2$ 的隐藏节点压缩后得到的树称为最简等效拓扑。
\end{definition}

在配电网中，等效拓扑是一个重要而务实的概念。若某个未量测接线点没有分支，它的位置不可由末端量测唯一确定；但总线路阻抗仍可辨识。对拓扑辨识而言，恢复可观测负荷之间的分支层级通常比恢复每个物理接头更可行。

\section{Additive phylogeny 思想}

Additive phylogeny 来自系统发育树重构。其基本想法是：若叶间距离是精确加性树度量，则可以递归地移除一个叶节点，先重构较小树，再把该叶节点接回适当位置。

对叶节点 $j$，其 limb length 可由
\begin{equation}
  \ell_j=\min_{i,k\in L\setminus\{j\}}
  \frac{d_{ij}+d_{jk}-d_{ik}}{2}
  \label{eq:limb_length}
\end{equation}
给出。在精确加性情形下，$\ell_j$ 是 $j$ 到其父分支点的边长。找到满足最小值的一对 $i,k$ 后，可在 $i$ 到 $k$ 的路径上按距离定位 $j$ 的接入点。

\begin{algorithm}[htbp]
  \caption{简化 Additive Phylogeny 递归框架}
  \label{alg:additive_phylogeny}
  \KwIn{叶集合 $L$ 上的精确加性距离矩阵 $D$}
  \KwOut{带隐藏节点的加权树 $\widetilde T$}
  \If{$|L|=2$}{
    返回一条连接两个叶节点的边，边权为 $d_{12}$\;
  }
  选择一个叶节点 $j$，用式\eqref{eq:limb_length}计算 $\ell_j$\;
  找到 $i,k$ 使 $d_{ik}=d_{ij}+d_{jk}-2\ell_j$\;
  从距离矩阵中暂时删除 $j$，递归重构 $L\setminus\{j\}$ 上的树\;
  在 $i$ 到 $k$ 的路径上距离 $i$ 为 $d_{ij}-\ell_j$ 的位置插入接入点 $h$\;
  添加边 $(h,j)$，权重为 $\ell_j$\;
  压缩或合并不必要的度为 $2$ 隐藏节点\;
  \KwRet{$\widetilde T$}
\end{algorithm}

\section{Neighbor joining 思想}

Neighbor joining 不要求事先指定根。它每次选择一对最可能共享最近父节点的叶子，合并为一个新簇。常用准则为
\begin{equation}
  Q(i,j)=(n-2)d_{ij}-\sum_{k\in L}d_{ik}-\sum_{k\in L}d_{jk},
\end{equation}
其中 $n=|L|$。选择 $Q(i,j)$ 最小的一对作为邻接叶，再估计它们到新隐藏节点的边长。Neighbor joining 对噪声有一定经验鲁棒性，因此常被用于近似树度量重构。

\section{Recursive grouping 思想}

Recursive grouping 更强调“由局部距离关系识别父子和兄弟簇”。在有根配电网中，如果根节点或台区总表可观测，可利用根到叶的距离辅助判断两个叶节点的最近公共祖先深度：
\begin{equation}
  h(\lca(i,j))=\frac{h(i)+h(j)-d_{ij}}{2}.
  \label{eq:lca_from_distance}
\end{equation}
若多个叶对的 $\lca$ 深度相同且大于与其他叶的公共深度，则它们可能共享同一隐藏分支点。

\begin{algorithm}[htbp]
  \caption{面向配电网的简化 recursive grouping}
  \label{alg:recursive_grouping}
  \KwIn{叶节点集合 $L$，根到叶深度估计 $h(i)$，叶间距离 $\widehat d_{ij}$，容差 $\tau$}
  \KwOut{最简等效隐藏树}
  对所有叶对计算 $\widehat h_{ij}=(h(i)+h(j)-\widehat d_{ij})/2$\;
  将 $\widehat h_{ij}$ 较大且彼此接近的叶对聚为候选 sibling group\;
  \ForEach{候选组 $C$}{
    若组内 $\widehat h_{ij}$ 的波动小于 $\tau$，创建隐藏父节点 $h_C$\;
    估计 $h_C$ 到每个叶 $i\in C$ 的边长为 $h(i)-\operatorname{median}_{j\in C}\widehat h_{ij}$\;
    把 $C$ 压缩为新 terminal $h_C$，更新到其他 terminal 的距离\;
  }
  重复分组直到不能产生新的稳定隐藏节点\;
  压缩所有度为 $2$ 的隐藏节点\;
  \KwRet{得到的等效树}
\end{algorithm}

该算法不是工业级实现，但展示了核心逻辑：从 LCA 深度或 quartet 关系中识别“共享分支”的叶节点簇，再递归向上合并。

\section{度为 2 隐藏节点的不可辨识性}

\Cref{prop:degree2_unidentifiable}说明，度为 $2$ 的隐藏节点不会改变 terminal 间距离的结构，只改变某条路径的拆分方式。隐藏树重构通常要求最简性：
\begin{equation}
  \deg(h)\ne 2,\qquad \forall h\in \widetilde H.
\end{equation}
这不是数学偷懒，而是可辨识性的边界。若需要恢复物理接头或中间杆塔，需要额外 GIS、施工台账或现场巡检信息。

\section{小规模数值例子}

\begin{example}[从四个叶节点恢复两个隐藏分支点]
\label{ex:hidden_tree_four_leaf}
设四个末端电表 $a,b,c,d$ 的估计距离为
\[
\begin{array}{c|cccc}
 & a&b&c&d\\ \hline
a&0&2&5&5\\
b&2&0&5&5\\
c&5&5&0&2\\
d&5&5&2&0
\end{array}
\]
由\cref{chap:four_point}的判别，$d_{ab}+d_{cd}=4$ 小于另外两项 $10$，故 quartet 分裂为 $ab|cd$。内部边长为 $3$。若设叶到各自隐藏父节点的边长均为 $1$，则得到隐藏树
\[
  a--u--v--c,\qquad b--u,\qquad d--v,
\]
其中 $w_{au}=w_{bu}=w_{cv}=w_{dv}=1$，$w_{uv}=3$。该树精确复现所有叶间距离。

若直接在 $\{a,b,c,d\}$ 上求 MST，会选 $(a,b)$、$(c,d)$ 以及任一条长度为 $5$ 的跨边。输出既没有 $u,v$，也无法表达 $a,b$ 与 $c,d$ 分别共享隐藏分支点。
\end{example}

\section{本章小结}

仅叶节点可观测时，拓扑辨识的目标应从“在观测节点上求生成树”转为“寻找诱导叶间距离的最简隐藏树”。Additive phylogeny 通过 limb length 递归接回叶节点，neighbor joining 通过邻接准则合并叶对，recursive grouping 利用 LCA 深度或局部分组识别共享分支。度为 $2$ 隐藏节点不可唯一恢复，因此最简等效拓扑是合理输出。

\section*{思考题}
\begin{enumerate}
  \item 对\cref{ex:hidden_tree_four_leaf}，如果跨簇距离从 $5$ 变为 $5.2,4.9,5.1,5.0$，应如何设定 quartet 容差？
  \item 为什么 neighbor joining 适合无根树，而配电网常常具有自然根节点？
  \item 若台区总表可观测，根到叶距离能给 hidden tree reconstruction 带来什么额外信息？
  \item 在实际低压台区中，哪些隐藏节点应被视为可辨识分支点，哪些应压缩进线路阻抗？
\end{enumerate}

\chapter{配电网中的阻抗距离}
\label{chap:impedance_distance}

\section*{本章目标}
本章把图论距离与配电网线性化潮流连接起来。我们从 LinDistFlow 近似出发，解释灵敏度矩阵 $R,X$ 的公共路径意义，构造电阻距离 $d^R$ 与电抗距离 $d^X$，并证明它们在树上是路径距离。因此，全节点可观测时可用 MST，叶节点可观测时则应结合\cref{chap:hidden_tree}的隐藏树重构。

\section{LinDistFlow 近似}

在径向配电网中，忽略二阶损耗项并在标称电压附近线性化，可得到电压幅值变化与功率注入变化之间的近似关系
\begin{equation}
  \Delta V \approx R\,\Delta P + X\,\Delta Q,
  \label{eq:lindistflow_matrix}
\end{equation}
其中 $\Delta V,\Delta P,\Delta Q\in\R^n$ 对应非根节点，$R$ 和 $X$ 是电压对有功、无功注入的灵敏度矩阵。不同文献中的常数因子可能因采用电压幅值、电压平方或 per-unit 标幺化而不同；本章关注其树路径结构。

\section{公共路径阻抗和}

设根节点为 $0$，非根节点集合为 $\{1,\ldots,n\}$。对边 $e$ 记电阻、电抗为 $r_e,x_e>0$。在线性化径向模型中，灵敏度矩阵具有公共路径形式：
\begin{align}
  R_{ij} &= \sum_{e\in \pathset_{0i}\cap \pathset_{0j}} r_e, \label{eq:R_common_path}\\
  X_{ij} &= \sum_{e\in \pathset_{0i}\cap \pathset_{0j}} x_e. \label{eq:X_common_path}
\end{align}
也就是说，$R_{ij}$ 是从根到 $i$ 与从根到 $j$ 两条路径的公共部分电阻和。公共部分正是从根到 $\lca(i,j)$ 的路径。

\begin{figure}[htbp]
  \centering
  \begin{tikzpicture}[
  every node/.style={font=\small},
  bus/.style={circle,draw,fill=blue!12,minimum size=7mm}
]
\node[bus] (n0) at (0,0) {$0$};
\node[bus] (n1) at (2,0) {$1$};
\node[bus] (n2) at (4,0.9) {$2$};
\node[bus] (n3) at (4,-0.9) {$3$};
\draw[very thick] (n0)--node[above] {$r=.10$} (n1);
\draw[very thick] (n1)--node[above] {$r=.08$} (n2);
\draw[very thick] (n1)--node[below] {$r=.12$} (n3);
\draw[dashed,blue!60] (n0) to[bend left=25] node[above] {公共路径} (n1);
\node[align=center,right=0.6cm of n3] {$R_{23}=r_{01}$\\$d^R_{23}=r_{12}+r_{13}$};
\end{tikzpicture}

  \caption{阻抗距离例子。$R_{ij}$ 等于根到 $\lca(i,j)$ 的公共路径电阻和。}
  \label{fig:impedance_tree}
\end{figure}

\section{阻抗距离的构造}

定义电阻距离与电抗距离：
\begin{align}
  d^R_{ij} &= R_{ii}+R_{jj}-2R_{ij}, \label{eq:resistance_distance}\\
  d^X_{ij} &= X_{ii}+X_{jj}-2X_{ij}. \label{eq:reactance_distance}
\end{align}

\begin{proposition}[阻抗距离是树上路径距离]
\label{prop:impedance_is_tree_distance}
若 $R_{ij}$ 与 $X_{ij}$ 满足公共路径形式\eqref{eq:R_common_path}--\eqref{eq:X_common_path}，则 $d^R_{ij}$ 等于 $i$ 到 $j$ 路径上的电阻和，$d^X_{ij}$ 等于 $i$ 到 $j$ 路径上的电抗和。
\end{proposition}

\begin{proof}
以电阻为例。$R_{ii}$ 是根到 $i$ 的路径电阻和，$R_{jj}$ 是根到 $j$ 的路径电阻和，$R_{ij}$ 是根到 $\lca(i,j)$ 的公共路径电阻和。因此
\begin{align}
  R_{ii}+R_{jj}-2R_{ij}
  &=
  \sum_{e\in \pathset_{0i}}r_e+
  \sum_{e\in \pathset_{0j}}r_e-
  2\sum_{e\in \pathset_{0i}\cap \pathset_{0j}}r_e\\
  &=
  \sum_{e\in \pathset_{ij}}r_e .
\end{align}
电抗情形完全相同。
\end{proof}

\begin{corollary}[阻抗距离与 MST/隐藏树]
\label{cor:impedance_mst_hidden}
若所有电气节点均可观测且可精确估计 $R$ 或 $X$，则以 $d^R$ 或 $d^X$ 为边权在完整节点集上求 MST 可恢复真实树。若只观测叶节点，则 $d^R,d^X$ 是叶间树距离，应使用隐藏树重构恢复最简等效拓扑。
\end{corollary}

\begin{proof}
由\cref{prop:impedance_is_tree_distance}，$d^R,d^X$ 是正权树路径距离。完整节点集情形应用\cref{thm:additive_mst_recovery}；叶节点情形则属于\cref{chap:hidden_tree}讨论的 terminal-only 重构问题。
\end{proof}

\section{三节点/四节点算例}

\begin{example}[四节点径向网的电阻距离]
考虑根节点 $0$，真实树边为 $(0,1)$、$(1,2)$、$(1,3)$，电阻分别为
\[
  r_{01}=0.10,\qquad r_{12}=0.08,\qquad r_{13}=0.12 .
\]
则
\[
R_{11}=0.10,\quad R_{22}=0.18,\quad R_{33}=0.22,\quad
R_{12}=0.10,\quad R_{13}=0.10,\quad R_{23}=0.10.
\]
由式\eqref{eq:resistance_distance}得
\[
d^R_{12}=0.08,\qquad d^R_{13}=0.12,\qquad d^R_{23}=0.20.
\]
若把根也纳入距离，$d^R_{01}=0.10$、$d^R_{02}=0.18$、$d^R_{03}=0.22$。在完整节点集 $\{0,1,2,3\}$ 上求 MST，会选择 $(1,2)$、$(0,1)$、$(1,3)$，恢复真实拓扑。
\end{example}

\section{物理意义}

阻抗距离的直观意义是“两个节点之间需要经过多少线路阻抗”。两个同一分支下游节点可能具有很大的源端公共路径，因此 $R_{ij}$ 大；但它们之间的路径距离 $d^R_{ij}$ 只计算从一个节点到另一个节点的非公共部分。直接把 $R_{ij}$ 当距离会把共享上游路径误读为远近关系；正确的距离构造必须使用 $R_{ii}+R_{jj}-2R_{ij}$。

\section{本章小结}

LinDistFlow 近似把电压变化写成 $\Delta V\approx R\Delta P+X\Delta Q$。在径向网络中，$R_{ij}$ 与 $X_{ij}$ 具有公共路径阻抗和结构。由此构造的 $d^R$ 与 $d^X$ 是树上路径距离：全节点可观测时可直接进入 MST 理论，叶节点可观测时进入隐藏树重构理论。这为从 $P/Q/V$ 曲线估计拓扑提供了物理桥梁。

\section*{思考题}
\begin{enumerate}
  \item 若某条线路电阻很小，MST margin 会如何变化？
  \item 为什么 $R_{ij}$ 本身不是距离？请检查它是否满足 $R_{ii}=0$。
  \item 在线性化模型中，三相不平衡和互阻抗会如何改变公共路径解释？
  \item 若只能估计 $R$ 而不能稳定估计 $X$，是否仍可辨识拓扑？需要哪些假设？
\end{enumerate}

\chapter{从 P/Q/V 曲线估计距离矩阵}
\label{chap:estimation}

\section*{本章目标}
本章说明如何从多时刻 $P/Q/V$ 曲线估计 LinDistFlow 灵敏度矩阵，并进一步构造阻抗距离矩阵。我们讨论岭回归、Lasso、Elastic Net、公共源端电压模态、未量测负荷导致的遗漏变量偏差，以及估计误差如何通过\cref{thm:noisy_margin_mst}影响 MST 恢复。

\section{多时刻线性模型}

对 $T$ 个时刻收集非根节点量测，构造差分或去均值矩阵
\begin{equation}
  \Delta V =
  \begin{bmatrix}
    \Delta V(1)^\top\\
    \vdots\\
    \Delta V(T)^\top
  \end{bmatrix},\quad
  \Delta P =
  \begin{bmatrix}
    \Delta P(1)^\top\\
    \vdots\\
    \Delta P(T)^\top
  \end{bmatrix},\quad
  \Delta Q =
  \begin{bmatrix}
    \Delta Q(1)^\top\\
    \vdots\\
    \Delta Q(T)^\top
  \end{bmatrix}.
\end{equation}
矩阵形式可写为
\begin{equation}
  \Delta V = \Delta P R^\top + \Delta Q X^\top + E,
  \label{eq:regression_matrix}
\end{equation}
等价于逐个电压节点 $i$ 回归
\begin{equation}
  \Delta V_i(t)=
  \sum_j R_{ij}\Delta P_j(t)+
  \sum_j X_{ij}\Delta Q_j(t)+E_i(t).
\end{equation}

若采用列向量约定，也常写成 $\Delta v=R\Delta p+X\Delta q+\epsilon$。关键不是转置位置，而是把电压响应与功率扰动联系起来。

\section{岭回归、Lasso 与 Elastic Net}

令设计矩阵为
\begin{equation}
  Z=[\Delta P\ \ \Delta Q]\in\R^{T\times 2n},
\end{equation}
对每个电压列 $y_i=\Delta V_{\cdot i}$ 估计系数
\begin{equation}
  \theta_i=[R_{i1},\ldots,R_{in},X_{i1},\ldots,X_{in}]^\top .
\end{equation}

\paragraph{岭回归}
\begin{equation}
  \widehat\theta_i^{\rm ridge}
  =\argmin_\theta \|y_i-Z\theta\|_2^2+\lambda\|\theta\|_2^2.
\end{equation}
岭回归适合特征强相关场景，可降低方差，但估计矩阵会被整体收缩。

\paragraph{Lasso}
\begin{equation}
  \widehat\theta_i^{\rm lasso}
  =\argmin_\theta \|y_i-Z\theta\|_2^2+\lambda\|\theta\|_1.
\end{equation}
Lasso 产生稀疏系数，但在配电网灵敏度矩阵中，$R_{ij}$ 本身并不稀疏；稀疏性更适合导纳矩阵或边集合。若直接对 $R$ 做强稀疏约束，可能破坏公共路径结构。

\paragraph{Elastic Net}
\begin{equation}
  \widehat\theta_i^{\rm en}
  =\argmin_\theta \|y_i-Z\theta\|_2^2
  +\lambda_1\|\theta\|_1+\lambda_2\|\theta\|_2^2.
\end{equation}
Elastic Net 在相关负荷下比纯 Lasso 更稳定，可作为工程折中。

\section{公共源端电压模态}

若源端电压 $V_0(t)$ 未完全稳定，则所有节点电压会共享共同扰动。观测模型更接近
\begin{equation}
  \Delta V(t)=\mathbf{1}\,\Delta V_0(t)+R\Delta P(t)+X\Delta Q(t)+E(t).
  \label{eq:source_mode}
\end{equation}
若忽略 $\Delta V_0(t)$，回归会把公共模态错误归因于负荷扰动，导致 $\widehat R,\widehat X$ 含有近似低秩偏差。

常见处理方式包括：
\begin{enumerate}[label=(\arabic*)]
  \item 若根节点电压可测，直接从所有节点电压中扣除 $\Delta V_0(t)$；
  \item 若不可测，用主成分或全体节点均值估计公共模态；
  \item 使用电压差 $V_i(t)-V_j(t)$ 消除共同源端项，但需重新推导回归形式；
  \item 在模型中显式加入截距、日周期项或公共因子。
\end{enumerate}

\section{未量测负荷与遗漏变量偏差}

若某些节点负荷未量测，真实模型为
\begin{equation}
  \Delta V = Z_o\Theta_o + Z_m\Theta_m + E,
\end{equation}
但估计时只使用观测设计矩阵 $Z_o$。若 $Z_m$ 与 $Z_o$ 相关，则遗漏项 $Z_m\Theta_m$ 与解释变量相关，普通最小二乘会产生偏差：
\begin{equation}
  \E[\widehat\Theta_o]-\Theta_o
  \approx (Z_o^\top Z_o)^{-1}Z_o^\top Z_m\Theta_m.
\end{equation}
在居民负荷同步性强的台区，这种偏差尤其明显。它会污染阻抗距离，进而缩小或翻转 MST 的 cut margin。

\section{估计误差与 MST margin}

假设估计得到 $\widehat R$，且
\begin{equation}
  \|\widehat R-R\|_{\max}\le \eta_R .
\end{equation}
由阻抗距离定义，
\begin{align}
  |\widehat d^R_{ij}-d^R_{ij}|
  &=
  |(\widehat R_{ii}-R_{ii})+(\widehat R_{jj}-R_{jj})
  -2(\widehat R_{ij}-R_{ij})|\\
  &\le 4\eta_R .
\end{align}
因此，若真实电阻距离的最小 cut margin 为 $\gamma_{\min}^R$，由\cref{thm:noisy_margin_mst}可得一个简单充分条件：
\begin{equation}
  4\eta_R < \frac{1}{2}\gamma_{\min}^R,
  \qquad\text{即}\qquad
  \eta_R<\frac{\gamma_{\min}^R}{8}.
\end{equation}
这个界通常保守，但清楚说明：统计估计精度必须与网络最小可分辨阻抗间隔匹配。

\section{Python 实现思路}

下面给出实现流程示意，不要求本书项目实际运行。

\begin{verbatim}
import numpy as np
from sklearn.linear_model import Ridge, ElasticNet

# P, Q, V: shape (T, n), already time-aligned
dP = P - P.mean(axis=0, keepdims=True)
dQ = Q - Q.mean(axis=0, keepdims=True)
dV = V - V.mean(axis=0, keepdims=True)

# optional: remove common voltage mode
dV = dV - dV.mean(axis=1, keepdims=True)

Z = np.hstack([dP, dQ])
coef = []
for i in range(n):
    model = Ridge(alpha=alpha, fit_intercept=False)
    model.fit(Z, dV[:, i])
    coef.append(model.coef_)
coef = np.asarray(coef)

R_hat = coef[:, :n]
X_hat = coef[:, n:]

dR = np.diag(R_hat)[:, None] + np.diag(R_hat)[None, :] - 2 * R_hat
dR = 0.5 * (dR + dR.T)  # enforce symmetry before MST
\end{verbatim}

实际工程还需处理缺失值、时间同步、异常点、三相相别、标幺化、温度和光伏反送等因素。

\section{本章小结}

从 $P/Q/V$ 曲线估计拓扑通常分两步：先用线性化潮流回归估计 $R,X$，再构造阻抗距离并进入 MST 或隐藏树重构。岭回归、Lasso、Elastic Net 各有偏差-方差权衡。公共源端电压和未量测负荷会造成系统性偏差；这种偏差最终体现为距离误差，并由 MST margin 决定是否改变拓扑输出。

\section*{思考题}
\begin{enumerate}
  \item 为什么直接回归 $V$ 而不做差分或去均值可能引入日周期偏差？
  \item 若 $\Delta P$ 与 $\Delta Q$ 高度相关，应如何解释 $R$ 与 $X$ 估计的不稳定性？
  \item 设计一个验证流程，判断公共源端电压模态是否已经被充分去除。
  \item 若估计误差主要集中在少数节点对，是否必须使用全局 $\|\cdot\|_{\max}$ 界？可以怎样改进 margin 分析？
\end{enumerate}

\chapter{通信数据与多视图图关系学习}
\label{chap:multiview}

\section*{本章目标}
本章讨论 PLC/HPLC 通信数据、GIS 台账与电气量测如何融合。重点观点是：通信图 $G_C$ 不等于电气图 $G_E$，但通信特征可以作为 noisy prior 帮助边概率估计。最后给出用最大生成树或有向生成树解码放射状拓扑的简化流程。

\section{PLC/HPLC 通信数据类型}

低压台区的 PLC/HPLC 系统可能提供如下特征：
\begin{enumerate}[label=(\arabic*)]
  \item RSSI：接收信号强度；
  \item SNR：信噪比，可包括整体 SNR 与子载波 SNR；
  \item RTT：往返时延；
  \item PER：包错误率或丢包率；
  \item hop：通信路由跳数；
  \item route：某时段路由父节点或中继路径；
  \item CIR：信道冲激响应；
  \item 频段特征：各子载波衰落、噪声谱、可用载波数量。
\end{enumerate}
这些特征与线路长度、阻抗、耦合方式、噪声源和中继策略相关，因此可能携带拓扑信息，但它们并不是电气边的直接观测。

\section{通信图不等于电气图}

定义通信图 $G_C=(V,E_C)$，若两个节点存在直接通信、路由邻接或信道质量超过阈值，则连接边。定义电气图 $G_E=(V,E_E)$，边表示实际导线或等效线路连接。一般有
\begin{equation}
  G_C\ne G_E.
\end{equation}
原因包括：
\begin{enumerate}[label=(\arabic*)]
  \item 通信信号可跨分支耦合，电气不相邻节点可能通信质量好；
  \item 中继路由由协议和瞬时信道决定，不必沿电气父子边；
  \item 电缆类型、接头、噪声源和相别会影响通信质量；
  \item HPLC 网络可能动态换路，通信边具有时间变动性。
\end{enumerate}

\section{通信距离作为 noisy prior}

通信特征可以转化为先验距离或边概率。例如
\begin{equation}
  \phi^{PLC}_{ij}=
  [\operatorname{RSSI}_{ij},\operatorname{SNR}_{ij},
  \operatorname{RTT}_{ij},\operatorname{PER}_{ij},
  \operatorname{hop}_{ij},\ldots].
\end{equation}
然后学习
\begin{equation}
  p_{ij}=
  \mathbb{P}\big((i,j)\in E_E
  \mid \phi^{V}_{ij},\phi^{PQ}_{ij},\phi^{PLC}_{ij},\phi^{GIS}_{ij}\big).
  \label{eq:edge_probability_multiview}
\end{equation}
若缺少标注拓扑，可用弱监督：GIS 边作为正样本候选，明显远距离或不同台区节点作为负样本候选，再用物理可行性筛选。

\section{多视图图}

可把不同信息源看成不同视图：
\begin{equation}
  G^V,\qquad G^{PQ},\qquad G^{PLC},\qquad G^{GIS}.
\end{equation}
其中 $G^V$ 来自电压响应或阻抗距离，$G^{PQ}$ 来自负荷相关与扰动共现，$G^{PLC}$ 来自通信特征，$G^{GIS}$ 来自台账候选。融合方式包括特征拼接、加权投票、贝叶斯先验、图神经网络或约束优化。本书建议入门阶段先采用透明模型：
\begin{equation}
  \operatorname{logit}(p_{ij})=
  \beta_0+\beta_V^\top\phi^V_{ij}
  +\beta_{PQ}^\top\phi^{PQ}_{ij}
  +\beta_{PLC}^\top\phi^{PLC}_{ij}
  +\beta_{GIS}^\top\phi^{GIS}_{ij}.
\end{equation}

\section{由边概率解码为放射状拓扑}

若网络被建模为无向放射状树，可求最大生成树：
\begin{equation}
  \widehat T=
  \argmax_{T\text{ spans }V}\sum_{(i,j)\in T}\log p_{ij}.
  \label{eq:maxst_probability}
\end{equation}
若已知根节点且希望恢复父子方向，可先求无向树再按根定向；若边概率本身有方向 $p_{i\to j}$，可求最大有向生成树或 arborescence。

\begin{algorithm}[htbp]
  \caption{多视图拓扑辨识简化流程}
  \label{alg:multiview_pipeline}
  \KwIn{$P/Q/V$ 曲线，PLC/HPLC 记录，GIS 候选边，根节点 $0$}
  \KwOut{放射状电气拓扑 $\widehat T$}
  从 $P/Q/V$ 估计阻抗距离与电气特征 $\phi^V,\phi^{PQ}$\;
  从通信记录提取 RSSI、SNR、RTT、PER、hop、route、CIR 特征 $\phi^{PLC}$\;
  从 GIS 提取候选关系、地理距离、同箱同相等特征 $\phi^{GIS}$\;
  用监督或弱监督模型估计每条候选边的 $p_{ij}$\;
  在候选图上求最大生成树；若需要方向，则以根节点 $0$ 定向\;
  用潮流残差、电压可行性和开关约束做后验校验\;
  \KwRet{$\widehat T$}
\end{algorithm}

\begin{figure}[htbp]
  \centering
  \begin{tikzpicture}[
  every node/.style={font=\small,align=center},
  box/.style={draw,rounded corners=2pt,minimum width=2.7cm,minimum height=0.8cm,fill=blue!8},
  fuse/.style={draw,rounded corners=2pt,minimum width=3.0cm,minimum height=0.9cm,fill=green!10},
  outbox/.style={draw,rounded corners=2pt,minimum width=3.0cm,minimum height=0.9cm,fill=orange!12},
  arr/.style={-Latex,thick}
]
\node[box] (v) at (0,2) {$P/Q/V$\\阻抗距离};
\node[box] (plc) at (0,0.65) {PLC/HPLC\\RSSI,SNR,route};
\node[box] (gis) at (0,-0.7) {GIS/台账\\候选边};
\node[fuse] (prob) at (4,0.65) {边概率模型\\$p_{ij}$};
\node[outbox] (mst) at (8,0.65) {最大生成树\\放射状解码};
\node[outbox] (check) at (8,-1.0) {潮流/约束\\后验校验};
\draw[arr] (v)--(prob);
\draw[arr] (plc)--(prob);
\draw[arr] (gis)--(prob);
\draw[arr] (prob)--(mst);
\draw[arr] (mst)--(check);
\end{tikzpicture}


  \caption{电气、通信、GIS 多视图融合并解码为放射状拓扑。}
  \label{fig:multiview_flow}
\end{figure}

\section{物理约束的作用}

多视图学习不能只追求分类精度。配电网拓扑还必须满足：
\begin{enumerate}[label=(\arabic*)]
  \item 每个非根节点有唯一路径连接根节点；
  \item 边集合与开关状态、馈线边界和相别约束一致；
  \item 在线性化或 AC 潮流下能解释观测电压；
  \item 不产生与 GIS 明显矛盾的长距离跳接，除非数据证据充分。
\end{enumerate}
这些约束可作为后处理，也可直接进入优化模型。

\section{本章小结}

PLC/HPLC 通信数据提供了拓扑辨识的重要辅助信息，但通信图不等于电气图。合理做法是把通信距离作为 noisy prior，与 $P/Q/V$ 估计的电气特征和 GIS 台账共同学习边概率，再用最大生成树或有向生成树解码出放射状拓扑。多视图方法的优势在于抗单一数据源失真，但必须保留物理约束和可解释校验。

\section*{思考题}
\begin{enumerate}
  \item 给出一个 HPLC 路由父节点不等于电气父节点的可能场景。
  \item 若 GIS 台账很旧，是否应把 GIS 候选边作为硬约束？为什么？
  \item 最大生成树解码与逐边阈值分类相比有什么优势？
  \item 如何设计无标注或弱标注数据下的边概率训练样本？
\end{enumerate}

\chapter{图拉普拉斯、关联矩阵与导纳矩阵}
\label{chap:laplacian}

\section*{本章目标}
本章介绍 incidence matrix、graph Laplacian、admittance matrix sparsity、reduced Laplacian、effective resistance 与 Kron reduction。这些对象把图论结构、网络方程和拓扑辨识联系起来，并与\cref{chap:impedance_distance}中的阻抗距离互相补充。

\section{关联矩阵}

给无向图 $G=(V,E)$ 任意指定边方向。节点-边关联矩阵 $B\in\R^{|V|\times |E|}$ 定义为
\begin{equation}
  B_{ve}=
  \begin{cases}
    1, & \text{若 }v\text{ 是边 }e\text{ 的尾端},\\
    -1,& \text{若 }v\text{ 是边 }e\text{ 的头端},\\
    0, & \text{否则}.
  \end{cases}
\end{equation}
方向只是代数约定，不改变无向图本身。

\begin{proposition}[树的关联矩阵秩]
\label{prop:incidence_rank}
若 $G$ 有 $n$ 个节点且连通，则 $\operatorname{rank}(B)=n-1$。特别地，若 $G$ 是树，则删去任意一行后的 reduced incidence matrix 为可逆方阵。
\end{proposition}

\begin{proof}
对任意边列，元素和为零，故 $\mathbf{1}^\top B=0$，秩至多为 $n-1$。若 $G$ 连通且 $B^\top x=0$，则对任意边 $(u,v)$ 有 $x_u=x_v$。连通性推出所有 $x_v$ 相等，故 $\operatorname{null}(B^\top)=\operatorname{span}\{\mathbf{1}\}$，从而 $\operatorname{rank}(B)=n-1$。树有 $n-1$ 条边，删去任意一行后得到 $(n-1)\times(n-1)$ 矩阵，秩为 $n-1$，所以可逆。
\end{proof}

\section{图拉普拉斯}

给边权 $g_e>0$，令 $W=\operatorname{diag}(g_e)$。加权图拉普拉斯定义为
\begin{equation}
  L=BWB^\top .
  \label{eq:laplacian_bwb}
\end{equation}
其元素满足
\begin{equation}
  L_{ij}=
  \begin{cases}
    \sum_{k:(i,k)\in E}g_{ik}, & i=j,\\
    -g_{ij}, & (i,j)\in E,\\
    0, & i\ne j,\ (i,j)\notin E.
  \end{cases}
\end{equation}
因此 $L$ 的非零非对角元直接反映图的稀疏邻接关系。

\begin{lemma}[拉普拉斯二次型]
\label{lem:laplacian_quadratic}
对任意 $x\in\R^{|V|}$，
\begin{equation}
  x^\top Lx=\sum_{(i,j)\in E}g_{ij}(x_i-x_j)^2.
\end{equation}
\end{lemma}

\begin{proof}
由 $L=BWB^\top$ 得 $x^\top Lx=(B^\top x)^\top W(B^\top x)$。边 $e=(i,j)$ 对应的 $(B^\top x)_e$ 等于 $x_i-x_j$ 或其相反数，平方后方向无关。
\end{proof}

\section{导纳矩阵稀疏性}

在配电网稳态电路模型中，节点导纳矩阵 $Y$ 的非对角元满足
\begin{equation}
  Y_{ij}=
  \begin{cases}
    -y_{ij},& (i,j)\in E,\\
    0,& (i,j)\notin E,
  \end{cases}
\end{equation}
对角元包含与节点相连的线路导纳和接地支路导纳。若忽略 shunt 或把其单独建模，$Y$ 与复权拉普拉斯非常接近。拓扑辨识可以看成恢复 $Y$ 的稀疏模式：
\begin{equation}
  (i,j)\in E \quad \Longleftrightarrow \quad Y_{ij}\ne 0,\quad i\ne j .
\end{equation}

这与\cref{chap:estimation}中的 $R,X$ 不同：$R,X$ 是路径公共阻抗矩阵，通常稠密；$Y$ 或拉普拉斯是边局部矩阵，通常稀疏。

\section{Reduced Laplacian}

配电网常把根节点 $0$ 作为参考电压。删除 $L$ 中根节点对应的行列，得到 reduced Laplacian $L_{\rm red}$。若图连通，则 $L_{\rm red}$ 正定。

\begin{proposition}[reduced Laplacian 正定]
\label{prop:reduced_laplacian_pd}
若 $G$ 连通且边权 $g_e>0$，则删除任意一个节点后的 $L_{\rm red}$ 是正定矩阵。
\end{proposition}

\begin{proof}
任取非零 $z\in\R^{n-1}$，把它补上被删除节点的值 $0$ 得到 $x\in\R^n$。若 $z\ne 0$，则 $x$ 不是常数向量。由连通性，存在边 $(i,j)$ 使 $x_i\ne x_j$。由\cref{lem:laplacian_quadratic}，$x^\top Lx>0$。而 $x^\top Lx=z^\top L_{\rm red}z$，故 $L_{\rm red}$ 正定。
\end{proof}

\section{Effective resistance}

电阻网络中的 effective resistance 定义为
\begin{equation}
  R^{\rm eff}_{ij}=(e_i-e_j)^\top L^+(e_i-e_j),
  \label{eq:effective_resistance}
\end{equation}
其中 $L^+$ 是 Moore--Penrose 伪逆。对一般网状图，有效电阻综合了所有并联路径；对树，它退化为唯一路径上的电阻和。

\begin{proposition}[树上的有效电阻等于路径电阻]
\label{prop:effective_tree_path}
若 $G$ 是一棵电阻树，边电导为 $g_e=1/r_e$，则任意 $i,j$ 的有效电阻等于 $i$ 到 $j$ 的路径电阻和。
\end{proposition}

\begin{proof}
在树中，从 $i$ 向 $j$ 注入单位电流时，电流只能沿唯一路径 $\pathset_{ij}$ 流动，路径外边电流为零。路径上每条边电压降为 $1\cdot r_e$，总电压差为 $\sum_{e\in\pathset_{ij}}r_e$。有效电阻定义为单位电流下的端电压差，故结论成立。
\end{proof}

\section{Kron reduction}

若内部节点 $H$ 未观测，观测节点为 $L$，可把导纳矩阵按节点分块：
\begin{equation}
  Y=
  \begin{bmatrix}
    Y_{LL} & Y_{LH}\\
    Y_{HL} & Y_{HH}
  \end{bmatrix}.
\end{equation}
消去隐藏节点电压得到 Kron reduced matrix
\begin{equation}
  Y^{\rm kron}=Y_{LL}-Y_{LH}Y_{HH}^{-1}Y_{HL}.
  \label{eq:kron_reduction}
\end{equation}
Kron reduction 保留观测节点的端口行为，但通常会使图变稠密：原来通过隐藏节点相连的多个叶节点会形成等效完全子图。这解释了为什么只看观测节点导纳或相似度时，不能简单把非零项解释为原始电气边。

\section{与拓扑辨识的关系}

本章对象在拓扑辨识中的作用可概括如下。
\begin{enumerate}[label=(\arabic*)]
  \item 关联矩阵描述边-节点约束，是潮流方程和径向约束的基础。
  \item 拉普拉斯和导纳矩阵的稀疏模式对应电气边，可用于参数估计。
  \item reduced Laplacian 给出参考节点下的可逆网络方程。
  \item effective resistance 在树上等于路径距离，可连接到 MST 与 four-point condition。
  \item Kron reduction 说明隐藏节点消去后会产生稠密等效耦合，提醒我们不要把观测端口图误认为原始拓扑。
\end{enumerate}

\section{本章小结}

关联矩阵、拉普拉斯和导纳矩阵提供了从图结构到电路方程的代数桥梁。拉普拉斯二次型刻画边差分能量，导纳矩阵稀疏性刻画直接电气邻接，effective resistance 在树上给出路径距离，Kron reduction 则解释了隐藏节点导致的端口等效稠密化。这些工具为物理约束拓扑辨识提供了严谨线性代数基础。

\section*{思考题}
\begin{enumerate}
  \item 为什么 $R,X$ 灵敏度矩阵通常稠密，而导纳矩阵通常稀疏？
  \item 对三节点链 $0$--$1$--$2$，写出关联矩阵和拉普拉斯矩阵。
  \item Kron reduction 后的非零边是否代表真实线路？结合隐藏星形网络说明。
  \item 在含环网络中，有效电阻为什么不再等于某条单一路径的电阻和？
\end{enumerate}

\chapter{总结与研究路线建议}
\label{chap:summary}

\section*{本章目标}
本章回顾全书主线：什么时候 MST 足够，什么时候必须使用 hidden tree reconstruction，什么时候需要物理约束，什么时候需要通信/GIS 先验。最后给出一个适合论文开题的研究问题与技术路线。

\section{何时 MST 足够}

MST 足够可靠的典型条件是：
\begin{enumerate}[label=(\arabic*)]
  \item 节点集接近完整电气节点集，隐藏分支点很少或已纳入模型；
  \item 边权是可加树路径距离，如由公共路径灵敏度构造的阻抗距离；
  \item 估计误差小于最小 cut margin 的一半；
  \item 网络在研究时段内确实放射状运行；
  \item 公共源端电压、未量测负荷和异常数据已被处理。
\end{enumerate}
在这些条件下，\cref{thm:additive_mst_recovery,thm:noisy_margin_mst}给出了清晰的理论保证。

\section{何时必须使用 hidden tree reconstruction}

若可观测节点主要是末端智能电表，而中间分支箱、接线点或开关柜未量测，就不应把叶节点 MST 解释为真实电气拓扑。此时应使用\cref{chap:four_point,chap:hidden_tree}中的 hidden tree reconstruction。输出应是最简等效拓扑，而不是含所有物理接头的施工图。

尤其要注意：
\begin{equation}
  \text{leaf-only distance matrix} \not\Rightarrow
  \text{MST recovers hidden branches}.
\end{equation}
能恢复隐藏分支，是因为叶间距离满足 four-point condition 并携带 quartet 信息，而不是因为 MST 会自动增加隐藏节点。

\section{何时需要物理约束}

当距离矩阵由统计相关性、机器学习特征或含噪回归得到时，必须引入物理约束：
\begin{enumerate}[label=(\arabic*)]
  \item LinDistFlow 或三相线性化潮流约束；
  \item 电压上下限与潮流残差校验；
  \item 导纳矩阵稀疏性或阻抗非负性；
  \item 放射状约束和候选开关状态约束；
  \item 源端公共模态和未量测负荷建模。
\end{enumerate}
物理约束的作用不是装饰模型，而是防止统计相似度偏离电气因果结构。

\section{何时需要通信/GIS 先验}

当 $P/Q/V$ 量测不足、扰动激励不够、负荷高度相关或局部 margin 很小时，通信与 GIS 先验尤其重要。PLC/HPLC 特征可提供边的弱证据，GIS 可限制候选边空间，开关台账可解释拓扑变化。但它们应作为 noisy prior，而不是直接替代电气拓扑。

\section{适合开题的研究问题}

一个自然且具有工程价值的研究问题是：

\begin{quote}
\textbf{融合 $P/Q/V$ 曲线分解与 PLC 通信特征的低压配电网多视图拓扑辨识。}
\end{quote}

可拆分为以下研究任务。
\begin{enumerate}[label=(\arabic*)]
  \item \textbf{电气距离估计}：从多时刻 $P/Q/V$ 曲线中分离公共源端模态，估计 $R,X$ 或端口阻抗距离。
  \item \textbf{隐藏树建模}：针对仅末端智能电表可观测场景，利用 four-point condition 或 recursive grouping 恢复最简等效分支拓扑。
  \item \textbf{通信先验融合}：提取 HPLC 的 RSSI、SNR、RTT、PER、route、CIR 和子载波特征，学习边概率先验。
  \item \textbf{放射状解码}：在候选图上用最大生成树、有向生成树或混合整数约束解码拓扑。
  \item \textbf{物理后验校验}：用 LinDistFlow/三相潮流残差、电压可行性和开关状态验证输出。
  \item \textbf{不确定性评估}：用 cut margin、bootstrap 或贝叶斯后验识别低置信度边，指导补充量测。
\end{enumerate}

\section{建议实验设计}

论文实验可分为四层：
\begin{enumerate}[label=(\arabic*)]
  \item \textbf{合成树实验}：控制边阻抗、隐藏节点比例、噪声强度，验证 MST margin 与 hidden tree reconstruction。
  \item \textbf{标准配电网算例}：在 IEEE 径向馈线或实际台区脱敏模型上生成 LinDistFlow 数据。
  \item \textbf{通信扰动实验}：模拟通信边与电气边不一致，检验多视图融合是否优于单一视图。
  \item \textbf{消融实验}：分别去掉公共模态校正、GIS 先验、PLC 特征和物理后验校验，观察拓扑边准确率、路径准确率和电压残差变化。
\end{enumerate}

推荐指标包括边准确率、父节点准确率、树编辑距离、叶间路径距离误差、潮流残差、低置信边召回率以及运行时间。

\section{本章小结}

本书的主线可以压缩为一句话：\textbf{MST 解决完整节点集上的加性距离树恢复，hidden tree reconstruction 解决仅叶节点距离下的隐藏分支恢复，物理约束与多视图先验解决实际量测偏差与信息不足。} 对研究生而言，关键不是记住某个算法名称，而是判断数据、节点集、距离语义与输出拓扑之间是否匹配。

\section*{思考题}
\begin{enumerate}
  \item 针对你手头的数据，判断它属于全节点可观测、仅叶节点可观测还是候选图内开关状态辨识。
  \item 设计一个实验验证“公共源端电压模态会误导相关性拓扑辨识”。
  \item 如果只能新增少量量测点，应优先布置在分支节点还是末端节点？请结合 hidden tree reconstruction 说明。
  \item 为上述开题题目写出三个可检验假设和对应评价指标。
\end{enumerate}


\backmatter
\nocite{*}
\bibliographystyle{plainnat}
\bibliography{refs}

\end{document}

